% cap_fundamentacao.tex
\chapter{Fundamentação Teórica}
\label{chap:fundamentacao_teorica}

Este capítulo apresenta os conceitos teóricos necessários para fundamentar a previsão horária multi-horizonte de material particulado fino (PM2,5) em séries temporais ambientais. O foco recai sobre dados faltantes, variáveis exógenas conhecidas no momento da previsão, redes neurais artificiais, modelos recorrentes, arquiteturas codificador-decodificador com atenção e modelos baseados em árvores com impulsionamento por gradiente (\textit{gradient boosting}).

\section{Material Particulado Fino (PM2,5) e Monitoramento da Qualidade do Ar}

O material particulado fino (PM2,5) é formado por partículas com diâmetro aerodinâmico de até 2,5 micrômetros. O diâmetro aerodinâmico é uma forma de expressar o comportamento da partícula no ar: partículas menores permanecem suspensas por mais tempo e podem ser inaladas com maior facilidade. Por seu tamanho reduzido, o PM2,5 pode atingir regiões profundas do sistema respiratório e está associado a efeitos respiratórios e cardiovasculares relevantes \cite{pope2006health,brook2010particulate}. As diretrizes globais da Organização Mundial da Saúde tratam PM2,5 e material particulado com diâmetro aerodinâmico de até 10 micrômetros (PM10) como poluentes prioritários para avaliação de qualidade do ar \cite{who2021air}.

PM10, por sua vez, representa partículas com diâmetro aerodinâmico de até 10 micrômetros. Embora seja uma fração mais grossa que PM2,5, sua concentração pode carregar informação sobre condições ambientais, emissão de partículas e episódios de maior concentração. Por isso, em problemas de previsão, PM10 pode ser usado como variável auxiliar quando está disponível de forma compatível com a série de PM2,5.

Em redes de monitoramento, sensores registram concentrações de poluentes em intervalos regulares e formam séries temporais. Uma série temporal de qualidade do ar, portanto, é uma sequência de medições ordenadas pelo tempo. A previsão de curto prazo busca estimar concentrações futuras a partir de medições recentes e de informações auxiliares conhecidas no momento da previsão. Em formulações horárias, é comum usar um bloco de observações passadas para prever um ou mais instantes futuros.

\section{Séries Temporais Multivariadas}

Uma série temporal é uma sequência de observações indexadas no tempo. Em qualidade do ar, a concentração de um poluente depende de autocorrelação, ciclos diários e semanais, condições meteorológicas, variações de emissão, mudanças de regime e eventos raros. Quando múltiplas variáveis são observadas simultaneamente, como PM2,5, PM10 e atributos de calendário, a tarefa passa a ser multivariada. Nesse caso, o modelo não observa apenas o histórico do alvo, mas também outras informações que podem ajudar a explicar sua variação.

A análise de séries temporais costuma separar padrões em componentes interpretáveis. \citeonline{hyndman2021forecasting} distinguem tendência, sazonalidade, ciclos e variações remanescentes como padrões úteis para escolher métodos de previsão. Em uma decomposição aditiva simples, uma observação pode ser representada por:

\[
    y_t = T_t + S_t + R_t,
\]

em que \(T_t\) representa tendência ou tendência-ciclo, \(S_t\) representa sazonalidade de período conhecido e \(R_t\) representa o restante da série, incluindo ruído, mudanças abruptas e eventos raros. Para séries horárias, pode haver mais de uma sazonalidade relevante, como ciclo diário, semanal e anual. Essa leitura fundamenta o uso de atributos temporais, termos periódicos e defasagens como formas de representar dependências temporais sem recorrer a informação futura.

Outro conceito central é autocorrelação, isto é, a associação entre a série e versões defasadas dela mesma. Para uma defasagem \(\ell\), a autocorrelação amostral pode ser escrita como:

\[
    \rho(\ell)=
    \frac{
        \sum_t (y_t-\bar{y})(y_{t-\ell}-\bar{y})
    }{
        \sum_t (y_t-\bar{y})^2
    }.
\]

Valores altos de \(\rho(1)\), \(\rho(24)\) ou \(\rho(168)\) indicam, respectivamente, persistência de curto prazo, repetição diária e repetição semanal. Em problemas de previsão, essa leitura ajuda a escolher janelas históricas, defasagens sazonais e termos cíclicos adequados ao período de amostragem da série.

Formalmente, seja $\mathbf{x}_t \in \mathbb{R}^{d}$ o vetor de atributos observado no instante $t$, $y_t \in \mathbb{R}$ a concentração de PM2,5 e $\mathbf{z}_{t+k}$ o vetor de atributos futuros conhecidos no instante de emissão da previsão, como calendário, termos periódicos ou previsões meteorológicas disponíveis operacionalmente. Uma janela de entrada é o bloco de valores passados entregue ao modelo; o horizonte é o conjunto de instantes futuros que se deseja prever. Para uma janela de entrada de tamanho $L$ e horizonte $H$, o problema pode ser escrito como:

\[
    \hat{\mathbf{y}}_{t+1:t+H}
    =
    f_{\theta}\left(
        \mathbf{x}_{t-L+1:t},
        \mathbf{z}_{t+1:t+H}
    \right),
\]

em que $f_{\theta}$ é o modelo treinado e $\hat{\mathbf{y}}_{t+1:t+H} = (\hat{y}_{t+1}, \ldots, \hat{y}_{t+H})$. Essa formulação explicita a diferença entre variáveis observadas apenas no passado e variáveis conhecidas para o horizonte futuro.

Em problemas de previsão, a ordem temporal dos dados deve ser preservada. Separações aleatórias, comuns em tarefas com registros independentes, podem introduzir vazamento de informação futura em séries temporais. Por isso, avaliações fora da amostra e divisões cronológicas são recomendadas para estimar desempenho em um cenário próximo ao uso real \cite{hyndman2021forecasting,tashman2000out,bergmeir2012crossvalidation}.

\subsection{Validação cruzada temporal e \textit{walk-forward}}

A validação cruzada tradicional divide os dados em partições e assume exemplos aproximadamente independentes. Essa suposição não vale diretamente para séries temporais, nas quais observações vizinhas compartilham autocorrelação, sazonalidade e tendência. Em previsão, uma avaliação coerente deve respeitar a direção do tempo: o modelo só pode ser treinado com dados anteriores ao período validado. Essa preocupação é central em testes fora da amostra para previsão \cite{tashman2000out} e motiva adaptações da validação cruzada para séries temporais \cite{bergmeir2012crossvalidation,hyndman2021forecasting}.

Uma adaptação comum é a validação cruzada temporal por janela expansiva, também chamada de \textit{walk-forward validation}. Nesse procedimento, cada partição usa uma janela de treino que termina antes do bloco de validação seguinte. A cada nova partição, a janela de treino se expande e a validação avança para um período futuro. Assim, a avaliação considera várias origens temporais de previsão sem permitir que dados futuros influenciem o treinamento. Para \(K\) partições, a métrica de seleção pode ser resumida pela média:

\[
    \mathcal{L}_{CV}
    =
    \frac{1}{K}\sum_{k=1}^{K}
    \mathcal{L}_{val}^{(k)}.
\]

O subscrito \(CV\) indica \textit{cross-validation}, isto é, validação cruzada. Em desenhos experimentais com seleção de hiperparâmetros, a validação cruzada temporal pode ser aplicada dentro de um bloco de desenvolvimento. Ela não substitui uma avaliação cronológica final separada: esse conjunto permanece reservado para estimar o desempenho fora da amostra depois da seleção. A Figura~\ref{fig:validacao_cruzada_temporal_walk_forward} ilustra uma janela expansiva com três partições e avaliação final reservada.

\begin{figure}[htbp]
    \centering
    \includegraphics[width=\linewidth]{validacao_cruzada_temporal_walk_forward.pdf}
    \caption[Validação cruzada temporal por janela expansiva]{Validação cruzada temporal por janela expansiva, ou \textit{walk-forward validation}. Cada partição treina com dados passados e valida no bloco temporal seguinte; a avaliação final permanece reservada para estimar desempenho fora da amostra. Elaboração própria com base em \citeonline{tashman2000out}, \citeonline{bergmeir2012crossvalidation} e \citeonline{hyndman2021forecasting}.}
    \label{fig:validacao_cruzada_temporal_walk_forward}
\end{figure}

\section{Previsão Multi-Horizonte}

Na previsão multi-horizonte, o objetivo é estimar mais de um instante futuro. Para uma entrada $\mathbf{x}_{t-L+1:t}$ e um horizonte de \(H\) instantes, o modelo deve produzir $\hat{\mathbf{y}}_{t+1:t+H}$. Revisões recentes de previsão com aprendizado profundo descrevem arquiteturas de codificador e decodificador tanto para previsão de horizonte único quanto para previsão multi-horizonte \cite{lim2021survey,benidis2022deep}. Três estratégias são especialmente relevantes:

\begin{itemize}
    \item \textbf{Previsão recursiva:} o modelo aprende uma função de horizonte único $\hat{y}_{t+1}=g_{\theta}(\mathbf{x}_{t-L+1:t})$ e realimenta suas próprias saídas para obter $\hat{y}_{t+2}, \ldots, \hat{y}_{t+H}$. A estratégia é simples, mas pode acumular erros ao longo do horizonte.
    \item \textbf{Previsão direta:} o modelo aprende diretamente $\hat{\mathbf{y}}_{t+1:t+H}=g_{\theta}(\mathbf{x}_{t-L+1:t})$. Isso reduz a realimentação de erro, mas exige aprender a estrutura conjunta do horizonte.
    \item \textbf{Codificador-decodificador:} uma rede codifica a janela de entrada e outra gera a sequência futura, possivelmente condicionada por variáveis conhecidas no horizonte. Essa formulação é natural para tarefas sequência-para-sequência.
\end{itemize}

Essas alternativas podem ser implementadas com modelos recorrentes, arquiteturas codificador-decodificador e modelos baseados em árvores com saída múltipla. Comparações desse tipo são necessárias porque a literatura de previsão não sustenta uma única arquitetura universalmente dominante; revisões de redes neurais recorrentes para séries temporais ressaltam que não existem soluções universais em previsão \cite{hewamalage2021rnn}.

\section{Dados Faltantes, Imputação e Vazamento Temporal}

Dados ambientais reais frequentemente contêm lacunas por falhas de sensor, manutenção, comunicação ou filtragem de leituras inconsistentes. Uma lacuna significa que, em determinado instante, a medição esperada não foi observada. A ausência de dados deve ser tratada com cuidado porque pode afetar tanto o treinamento quanto a avaliação. A literatura de dados faltantes diferencia o valor não observado do valor realmente medido; tratá-los como equivalentes pode distorcer inferências e métricas \cite{rubin1976inference}.

Imputação é o processo de preencher lacunas com estimativas. Métodos simples, como interpolação linear, preservam a estrutura temporal local e são transparentes; métodos multivariados, como Multivariate Imputation by Chained Equations (MICE) e MissForest, exploram relações entre variáveis \cite{vanbuuren2011mice,stekhoven2012missforest}. Métodos de imputação temporal, como Bidirectional Recurrent Imputation for Time Series (BRITS) e Self-Attention-based Imputation for Time Series (SAITS), também usam máscaras de observação para distinguir valores disponíveis de valores ausentes \cite{cao2018brits,du2023saits}. Neste contexto, máscara significa um indicador que informa se um valor foi originalmente observado ou se foi preenchido por algum método de imputação. Independentemente do método escolhido, essa distinção é importante: valores preenchidos não devem ser tratados automaticamente como observações reais. Por isso, protocolos de previsão com lacunas precisam preservar informação sobre ausência e explicitar como a imputação afeta treinamento e avaliação.

Outro risco é o vazamento temporal. A imputação, a normalização e a criação de atributos precisam ser ajustadas sem usar estatísticas futuras indisponíveis no instante de previsão. Em um protocolo cronológico, esses procedimentos devem ser definidos com informação anterior ao período avaliado.

\section{Janelas e Variáveis Exógenas}

Modelos supervisionados de séries temporais costumam converter uma sequência contínua em pares de entrada e saída por janelas deslizantes. Cada amostra é formada por uma janela histórica de \(L\) instantes e por um alvo futuro com \(H\) instantes. Em termos simples, o modelo recebe um trecho do passado e aprende a produzir um trecho do futuro. A sobreposição entre janelas aumenta o número de exemplos, mas exige cuidado para manter a divisão temporal e evitar que uma janela de treino alcance o período de validação ou teste.

Defasagens explícitas e variáveis exógenas podem aumentar a informação disponível para o modelo. Defasagem é o uso de valores passados de uma variável, como a concentração observada uma hora ou vinte e quatro horas antes. Variável exógena é uma informação externa ao alvo principal, como PM10, temperatura, vento ou atributos de calendário. No entanto, uma variável exógena só pode entrar no decodificador se estiver disponível no momento em que a previsão seria feita. Atributos de calendário, termos periódicos e previsões meteorológicas publicadas antes da emissão da previsão são exemplos de informação potencialmente disponível. Já variáveis observadas apenas no futuro, como medições futuras de outro poluente, não devem ser usadas diretamente.

\section{Redes Neurais e Aprendizado Profundo}
\label{sec:redes_neurais_deep_learning}

Redes neurais artificiais são modelos paramétricos compostos por unidades simples conectadas por pesos. A inspiração histórica vem de modelos computacionais do neurônio biológico, como o neurônio lógico de \citeonline{mcculloch1943logical} e o perceptron de \citeonline{rosenblatt1958perceptron}. Essa relação, porém, deve ser entendida como uma analogia didática. Redes neurais usadas em aprendizado de máquina não simulam de forma completa a fisiologia de neurônios reais; elas usam uma abstração matemática capaz de combinar entradas, ajustar pesos e produzir saídas.

Em um neurônio biológico, dendritos recebem sinais, o corpo celular integra parte dessas informações e o axônio transmite sinais para outras células. Em um neurônio artificial, a operação análoga é representada por entradas numéricas, pesos, viés e uma função de ativação. A Figura~\ref{fig:neuronio_biologico_artificial} apresenta essa relação de forma visual.

\begin{figure}[htbp]
    \centering
    \begin{subfigure}[t]{0.48\linewidth}
        \centering
        \includegraphics[width=\linewidth,trim=100 500 320 520,clip]{neuronio-biologico-wikimedia-ptbr.png}
        \caption{Neurônio biológico.}
        \label{fig:neuronio_biologico}
    \end{subfigure}
    \hfill
    \begin{subfigure}[t]{0.48\linewidth}
        \centering
        \includegraphics[width=0.86\linewidth]{neuronio-artificial-perceptron-wikimedia.png}
        \caption{Neurônio artificial.}
        \label{fig:neuronio_artificial_perceptron}
    \end{subfigure}
    \caption[Comparação visual entre neurônio biológico e neurônio artificial]{Comparação visual entre neurônio biológico e neurônio artificial. A imagem do neurônio biológico possui rótulos em português do Brasil e foi obtida de \citeonline{vasconcelos2008neuronPt}; a imagem do perceptron explicita entradas \(x_i\), pesos \(w_i\), soma ponderada, função de ativação \(\varphi\) e saída \(o\), conforme \citeonline{thoma2014perceptronUnit}.}
    \label{fig:neuronio_biologico_artificial}
\end{figure}

Formalmente, para um vetor de entrada $\mathbf{x}=(x_1,\ldots,x_d)$, um neurônio artificial calcula uma soma ponderada e aplica uma função de ativação:

\[
    a = \sum_{j=1}^{d} w_jx_j + b,
    \quad
    h = \phi(a).
\]

Nessa expressão, $w_j$ são pesos, $b$ é o viés, $a$ é a ativação antes da não linearidade e $\phi$ é uma função de ativação, como sigmoide, tangente hiperbólica ou unidade linear retificada (\textit{Rectified Linear Unit}, ReLU). Na Figura~\ref{fig:neuronio_artificial_perceptron}, o viés é representado por um peso \(w_0\) associado a uma entrada fixa \(x_0=1\), notação equivalente à adição explícita do termo \(b\). Os pesos e vieses são parâmetros aprendidos a partir dos dados. A função de ativação é importante porque introduz não linearidade; sem ela, uma composição de camadas lineares continuaria equivalente a uma única transformação linear.

Uma rede neural organiza vários neurônios em camadas. A saída de uma camada alimenta a camada seguinte, permitindo que o modelo componha transformações simples em representações mais complexas. De forma resumida, uma camada pode ser escrita como:

\[
    \mathbf{h}^{(\ell)}
    =
    \phi\left(
        W^{(\ell)}\mathbf{h}^{(\ell-1)}
        +
        \mathbf{b}^{(\ell)}
    \right),
\]

em que \(W^{(\ell)}\) e \(\mathbf{b}^{(\ell)}\) são os parâmetros da camada \(\ell\). O treinamento ajusta esses parâmetros para minimizar uma função de perda nos exemplos observados. Em tarefas de regressão, essa perda mede a distância entre os valores previstos e observados.

Aprendizado profundo, ou \textit{deep learning}, refere-se ao treinamento de redes neurais com múltiplas camadas e representações sucessivas dos dados \cite{goodfellow2016deep,lecun2015deep}. Em séries temporais, essa ideia é útil porque o modelo pode aprender relações não lineares entre medições recentes, variáveis exógenas, padrões cíclicos e o horizonte previsto. As arquiteturas recorrentes e codificador-decodificador discutidas nas seções seguintes são casos específicos de redes profundas voltadas a sequências: elas compartilham parâmetros ao longo do tempo, mantêm estados internos e produzem previsões para um ou mais horizontes futuros.

\section{Redes Neurais Recorrentes e Long Short-Term Memory (LSTM)}
\label{sec:rnn_lstm}

Redes neurais recorrentes (Recurrent Neural Networks, RNNs) são projetadas para dados sequenciais, pois atualizam um estado interno conforme processam cada instante da sequência. Em RNNs simples, o treinamento pode sofrer com desaparecimento ou explosão de gradientes, dificultando o aprendizado de dependências mais longas. A arquitetura Long Short-Term Memory (LSTM) introduz uma célula de memória e portões de entrada, esquecimento e saída para controlar o fluxo de informação \cite{hochreiter1997long}.

Uma célula LSTM pode ser descrita pelas seguintes equações, para entrada $\mathbf{x}_t$, estado oculto anterior $\mathbf{h}_{t-1}$ e memória anterior $\mathbf{c}_{t-1}$:

\[
\begin{aligned}
    \mathbf{i}_t &= \sigma(W_i\mathbf{x}_t + U_i\mathbf{h}_{t-1} + \mathbf{b}_i),\\
    \mathbf{f}_t &= \sigma(W_f\mathbf{x}_t + U_f\mathbf{h}_{t-1} + \mathbf{b}_f),\\
    \mathbf{o}_t &= \sigma(W_o\mathbf{x}_t + U_o\mathbf{h}_{t-1} + \mathbf{b}_o),\\
    \tilde{\mathbf{c}}_t &= \tanh(W_c\mathbf{x}_t + U_c\mathbf{h}_{t-1} + \mathbf{b}_c),\\
    \mathbf{c}_t &= \mathbf{f}_t \odot \mathbf{c}_{t-1} + \mathbf{i}_t \odot \tilde{\mathbf{c}}_t,\\
    \mathbf{h}_t &= \mathbf{o}_t \odot \tanh(\mathbf{c}_t).
\end{aligned}
\]

O portão de esquecimento $\mathbf{f}_t$ controla quanto da memória anterior é preservado; o portão de entrada $\mathbf{i}_t$ controla quanto da nova informação candidata é incorporado; e o portão de saída $\mathbf{o}_t$ controla a parte da memória exposta no estado oculto. A Figura~\ref{fig:lstm_propria} resume essa estrutura.

\begin{figure}[htbp]
    \centering
    \includegraphics[width=0.95\textwidth]{diagrama_lstm_proprio.pdf}
    \caption{Diagrama conceitual de uma célula LSTM. Elaboração própria com base em \citeonline{hochreiter1997long}.}
    \label{fig:lstm_propria}
\end{figure}

Em previsão multi-horizonte, a LSTM pode ser usada em estratégias recursivas ou diretas. Na estratégia recursiva, o modelo aprende a avançar ao longo do horizonte com realimentação. Na estratégia direta, o estado final da janela de entrada é projetado para todo o horizonte futuro. A comparação entre essas estratégias permite separar o efeito da arquitetura recorrente do efeito da forma de gerar o horizonte.

\section{Sequence-to-Sequence (Seq2Seq) e Atenção}
\label{sec:encoder_decoder_attention}

Modelos \textit{Sequence-to-Sequence} (Seq2Seq) foram propostos para mapear uma sequência de entrada em uma sequência de saída. \citeonline{sutskever2014sequence} descrevem a proposta como uma abordagem ``\textit{end-to-end}'' para aprendizado de sequências, em que um codificador transforma a entrada em representação latente e um decodificador gera a sequência de saída. Em séries temporais, essa estrutura é útil porque o decodificador pode produzir um horizonte completo e receber variáveis conhecidas para cada horizonte futuro.

Sejam $\mathbf{h}_j$ os estados do codificador para os \(L\) instantes da janela de entrada. Uma formulação simplificada do codificador é:

\[
    \mathbf{h}_j = Enc_{\theta}(\mathbf{x}_{t-L+j}, \mathbf{h}_{j-1}), \quad j=1,\ldots,L.
\]

Sem atenção, o decodificador depende principalmente de um vetor de contexto, frequentemente associado ao último estado do codificador. A atenção reduz esse gargalo. O próprio trabalho de \citeonline{bahdanau2014neural} caracteriza o ``\textit{fixed-length vector}'' como uma limitação em tarefas sequência-para-sequência. Com atenção aditiva, para cada horizonte futuro \(k\), calculam-se escores, pesos normalizados e um contexto:

\[
\begin{aligned}
    e_{k,j} &= \mathbf{v}_a^{\top}\tanh(W_s\mathbf{s}_{k-1} + W_h\mathbf{h}_{j} + \mathbf{b}_a),\\
    \alpha_{k,j} &= \frac{\exp(e_{k,j})}{\sum_{\ell=1}^{L}\exp(e_{k,\ell})},\\
    \mathbf{c}_k &= \sum_{j=1}^{L}\alpha_{k,j}\mathbf{h}_j.
\end{aligned}
\]

O decodificador atualiza seu estado e produz a previsão:

\[
\begin{aligned}
    \mathbf{s}_k &= Dec_{\theta}(\mathbf{s}_{k-1}, \tilde{y}_{t+k-1}, \mathbf{z}_{t+k}, \mathbf{c}_k),\\
    \hat{y}_{t+k} &= g_{\theta}(\mathbf{s}_k, \mathbf{c}_k, \mathbf{z}_{t+k}).
\end{aligned}
\]

Nessas expressões, \(W_s\), \(W_h\), \(\mathbf{v}_a\) e \(\mathbf{b}_a\) são parâmetros aprendidos do mecanismo de atenção, enquanto \(\mathbf{s}_k\) representa o estado do decodificador no horizonte \(k\).

Nesse caso, $\tilde{y}_{t+k-1}$ representa o valor real anterior durante \textit{teacher forcing} ou a previsão anterior durante inferência. Variações de atenção global e local também foram propostas, incluindo mecanismos do tipo Luong \cite{luong2015effective}. Para séries temporais multivariadas, arquiteturas como Dual-stage Attention-based Recurrent Neural Network (DA-RNN) combinam atenção sobre variáveis de entrada e atenção temporal \cite{qin2017darnn}.

A utilidade de Seq2Seq em previsão não se apoia na ideia de que essa família seja sempre superior. A justificativa é mais específica: tarefas multi-horizonte podem se beneficiar de um decodificador condicionado por atributos futuros conhecidos e trabalhos de previsão adotam estruturas recorrentes sequência-para-sequência para prever uma sequência de horizontes futuros. Por exemplo, \citeonline{wen2017mqrnn} exploram a natureza temporal de ``\textit{Sequence-to-Sequence Neural Networks}'' em previsão multi-horizonte, enquanto DeepAR e Temporal Fusion Transformer exemplificam linhas modernas de previsão profunda com decodificadores autoregressivos, covariáveis e previsão multi-horizonte \cite{salinas2020deepar,lim2021tft}. A Figura~\ref{fig:seq2seq_attention_propria} sintetiza uma formulação conceitual de Seq2Seq com atenção.

\begin{figure}[htbp]
    \centering
    \includegraphics[width=0.95\textwidth]{diagrama_seq2seq_attention_proprio.pdf}
    \caption[Seq2Seq com atenção]{Diagrama conceitual de Seq2Seq com atenção para previsão multi-horizonte. Elaboração própria com base em \citeonline{sutskever2014sequence}, \citeonline{bahdanau2014neural}, \citeonline{luong2015effective} e \citeonline{qin2017darnn}.}
    \label{fig:seq2seq_attention_propria}
\end{figure}

Pesos de atenção não devem ser interpretados automaticamente como explicação causal \cite{jain2019attention} nem como garantia de superioridade. Eles podem ser pouco informativos quando ficam quase uniformes ou quando não produzem ganho mensurável em avaliação fora da amostra.

\section{\textit{Teacher Forcing} e \textit{Scheduled Sampling}}

Em modelos com decodificador autoregressivo, o treinamento pode usar \textit{teacher forcing}: em vez de alimentar a previsão anterior do próprio modelo, usa-se o valor observado anterior como entrada do próximo instante. Isso acelera o treinamento, mas cria diferença entre treino e inferência, já que o valor real futuro não está disponível no momento da previsão. \citeonline{bengio2015scheduled} descrevem esse problema como discrepância entre treinamento e inferência, pois o item anterior verdadeiro está disponível no treino, mas é substituído por uma saída gerada pelo próprio modelo na inferência.

Seja \(p_{TF}\) a probabilidade de usar o valor real anterior no decodificador. Em cada horizonte futuro \(k\), pode-se sortear \(b_{i,k}\sim\text{Bernoulli}(p_{TF})\) e definir a entrada autoregressiva como:

\[
    \tilde{y}_{t+k-1}
    =
    b_{i,k}y_{t+k-1}
    +
    (1-b_{i,k})\hat{y}_{t+k-1}.
\]

Quando \(p_{TF}=1\), o modelo opera em \textit{teacher forcing} puro; quando \(p_{TF}=0\), ele treina de modo equivalente à inferência autoregressiva, sempre realimentando suas próprias previsões. O \textit{scheduled sampling} clássico propõe variar \(p_{TF}\) ao longo do treinamento, geralmente reduzindo essa probabilidade para deslocar gradualmente o decodificador de um regime totalmente guiado para um regime menos guiado e mais próximo da inferência \cite{bengio2015scheduled}.

Para previsão de séries temporais contínuas, entretanto, há uma motivação adicional para testar a direção oposta. \citeonline{teutsch2022flipped} argumentam que o \textit{scheduled sampling} clássico, herdado de tarefas de processamento de linguagem natural (Natural Language Processing, NLP), pode não ser a estratégia mais adequada para previsão e propõem currículos de \textit{Increasing Teacher Forcing}, denominados \textit{Flipped Classroom}. Nessa família, o treinamento começa mais próximo de \textit{free running}, com baixo \(p_{TF}\), e aumenta gradualmente a presença de \textit{teacher forcing}. Em experimentos com sistemas dinâmicos, os autores relataram ganhos de aproximadamente 15\% a 80\% na raiz do erro quadrático médio normalizado (Normalized Root Mean Squared Error, NRMSE) em previsão multi-horizonte em relação à melhor configuração de referência de \textit{teacher forcing} ou \textit{free running}, embora limitem a conclusão ao escopo experimental avaliado.

\clearpage

\begin{figure}[htbp]
    \centering
    \includegraphics[width=\linewidth]{teacher_forcing_scheduled_sampling.pdf}
    \caption[\textit{Teacher forcing}, \textit{scheduled sampling} e agenda crescente]{Comparação entre \textit{teacher forcing}, inferência autoregressiva, \textit{scheduled sampling} clássico e agenda crescente de \textit{teacher forcing}. Elaboração própria com base em \citeonline{bengio2015scheduled} e \citeonline{teutsch2022flipped}.}
    \label{fig:teacher_forcing_scheduled_sampling}
\end{figure}

A notação da Figura \ref{fig:teacher_forcing_scheduled_sampling} é importante: \(p_{TF}\) representa a probabilidade de usar o valor real anterior. Assim, uma agenda decrescente de \(p_{TF}\) aumenta o uso das próprias previsões do modelo e corresponde à intuição original do \textit{scheduled sampling}; já uma agenda crescente de \(p_{TF}\) aumenta a dependência de \textit{teacher forcing} e se aproxima da hipótese de \textit{Increasing Teacher Forcing}. Essa distinção evita confundir a direção da agenda de treinamento com a interpretação do regime autoregressivo.

\clearpage

\section{Impulsionamento por Gradiente e Extreme Gradient Boosting (XGBoost)}

Modelos de árvores com impulsionamento (\textit{boosting}) constroem um preditor aditivo por meio de várias árvores fracas, ajustadas sequencialmente para corrigir erros residuais. O \textit{Extreme Gradient Boosting} (XGBoost) é uma implementação eficiente de impulsionamento por gradiente (\textit{gradient boosting}) com regularização, paralelização e recursos práticos para matrizes de atributos \cite{chen2016xgboost}. Neste trabalho, essa matriz representa janelas temporais transformadas em atributos numéricos para uma tarefa supervisionada. Por lidar bem com atributos defasados e variáveis derivadas, o XGBoost é frequentemente usado como modelo de referência para previsão de séries temporais transformadas nesse formato.

No impulsionamento aditivo, a predição após $K$ árvores pode ser escrita como:

\[
    \hat{y}_i^{(K)} = \sum_{k=1}^{K} f_k(\mathbf{u}_i), \quad f_k \in \mathcal{F},
\]

em que $\mathbf{u}_i$ é o vetor de atributos derivado da janela temporal. O objetivo regularizado do XGBoost pode ser representado por:

\[
    \mathcal{L}^{(K)}
    =
    \sum_i \ell(y_i, \hat{y}_i^{(K)})
    +
    \sum_{k=1}^{K}\Omega(f_k),
    \quad
    \Omega(f_k)=\gamma T_k + \frac{1}{2}\lambda\lVert\mathbf{w}_k\rVert_2^2.
\]

Nessa expressão, $T_k$ é o número de folhas da árvore e $\mathbf{w}_k$ são os pesos das folhas. A regularização penaliza árvores excessivamente complexas e ajuda a controlar sobreajuste. Em problemas multi-horizonte, o XGBoost pode ser combinado com estratégias de saída múltipla para prever todo o horizonte futuro de uma vez; a Figura~\ref{fig:xgboost_multioutput_propria} ilustra essa conversão por janelas.

\begin{figure}[htbp]
    \centering
    \includegraphics[width=0.90\textwidth]{diagrama_xgboost_multioutput_proprio.pdf}
    \caption{Conversão de uma janela temporal em matriz de atributos para XGBoost com saída múltipla. Elaboração própria com base em \citeonline{chen2016xgboost}.}
    \label{fig:xgboost_multioutput_propria}
\end{figure}

Em aplicações com dados faltantes, a conversão por janelas precisa manter consistência entre as entradas, o horizonte previsto e os critérios de inclusão de alvos observados. Essa decisão afeta a comparabilidade entre modelos baseados em árvores e modelos neurais quando há máscaras ou perdas definidas ponto a ponto.

\section{Treinamento, Função de Perda e Otimização de Hiperparâmetros}
\label{sec:otimizacao_hpo}

Parâmetros são aprendidos pelo modelo durante o treinamento; hiperparâmetros, como taxa de aprendizado, número de camadas, tamanho do estado oculto e \textit{dropout}, são definidos antes ou buscados por procedimento externo. Essa busca é chamada de otimização de hiperparâmetros, indicada neste trabalho pela sigla HPO, do inglês \textit{Hyperparameter Optimization}. Redes neurais são frequentemente treinadas com otimizadores baseados em gradiente, como Adam e AdamW \cite{kingma2014adam,loshchilov2019decoupled}, e podem usar monitoramento em validação para seleção de ponto de controle.

A taxa de aprendizado controla o tamanho do passo aplicado aos parâmetros a cada atualização por gradiente. Ela pode ser mantida constante ou seguir uma agenda ao longo das épocas. Uma estratégia comum combina aquecimento linear (\textit{linear warm-up}) e decaimento cosseno: o aquecimento linear aparece em protocolos de treinamento com lotes grandes \cite{goyal2017accurate}, enquanto agendas cossenoidais são discutidas no contexto de reinícios quentes \cite{loshchilov2017sgdr}. Com \(E_w\) épocas de aquecimento, \(E\) épocas máximas, taxa máxima \(\eta_{\max}\) e taxa mínima \(\eta_{\min}\), a forma idealizada é:

\[
    \eta_t =
    \begin{cases}
        \eta_{\max}\left(10^{-5} + (1-10^{-5})\frac{t}{E_w}\right), & 0 \leq t \leq E_w,\\
        \eta_{\min} + \frac{1}{2}(\eta_{\max}-\eta_{\min})\left[1+\cos\left(\pi\frac{t-E_w}{E-E_w}\right)\right], & E_w < t \leq E.
    \end{cases}
\]

Assim, as primeiras épocas usam passos pequenos e aumentam gradualmente a taxa até \(\eta_{\max}\). Depois do aquecimento, a taxa decai suavemente até \(\eta_{\min}\). Essa variação ocorre por época durante o treinamento. A etapa de validação, por sua vez, não aplica retropropagação nem atualiza os pesos; ela avalia o modelo no fim das épocas e fornece métricas para parada antecipada, escolha de ponto de controle e poda de tentativas na HPO. Portanto, a validação influencia a seleção e a interrupção de treinamentos, mas não ajusta diretamente a taxa de aprendizado como ocorreria em uma agenda dependente da métrica de validação.

Também são usadas funções de perda robustas ou ponderadas. A perda de Huber combina comportamento quadrático para erros pequenos e linear para erros grandes, reduzindo a sensibilidade a outliers em comparação com o erro quadrático puro \cite{huber1964robust}. Para erro $e=y-\hat{y}$, ela é definida por:

\[
    \ell_{\delta}(e)=
    \begin{cases}
        \frac{1}{2}e^2, & \text{se } |e| \leq \delta,\\
        \delta\left(|e|-\frac{1}{2}\delta\right), & \text{caso contrário}.
    \end{cases}
\]

Quando um conjunto de dados contém valores imputados no alvo, a função de perda e as métricas de avaliação podem usar uma máscara \(m_{i,k}\) que vale 1 quando o alvo do exemplo \(i\) e horizonte \(k\) foi observado e 0 quando foi imputado. Essa escolha segue a distinção entre valor observado e valor ausente ou estimado, além do uso de máscaras de observação em métodos de imputação temporal \cite{rubin1976inference,cao2018brits,du2023saits}:

\[
    \mathcal{J}(\theta)
    =
    \frac{\sum_{i,k}m_{i,k}\ell(y_{i,k}, \hat{y}_{i,k})}
    {\sum_{i,k}m_{i,k}}.
\]

Uma perda absoluta ponderada altera explicitamente o objetivo para atribuir maior peso a determinados regimes de concentração ou tipos de erro:

\[
    \mathcal{J}_{w}(\theta)
    =
    \frac{\sum_{i,k}m_{i,k}w_{i,k}|y_{i,k}-\hat{y}_{i,k}|}
    {\sum_{i,k}m_{i,k}w_{i,k}}.
\]

Essa formulação não é regularização L1 dos parâmetros; é uma perda absoluta ponderada no espaço das previsões. Por isso, eventuais mudanças de comportamento em picos ou cauda devem ser interpretadas como consequência de uma mudança de objetivo, e não como simples ganho arquitetural.

O Optuna é uma ferramenta para HPO com definição dinâmica de espaços de busca, estratégias de amostragem e poda de tentativas pouco promissoras \cite{akiba2019optuna}. Formalmente, seja $\lambda \in \Lambda$ um vetor de hiperparâmetros no espaço de busca $\Lambda$ e seja $A(D_{treino};\lambda)$ o algoritmo de treinamento com essa configuração. Em uma validação simples, o objetivo da HPO seria selecionar a configuração com menor erro de validação:

\[
    \lambda^{*}
    =
    \arg\min_{\lambda \in \Lambda}
    \mathcal{L}_{val}\left(A(D_{treino};\lambda), D_{val}\right).
\]

Em séries temporais, esse termo de validação pode ser substituído pela média das partições temporais, isto é, pela métrica \(\mathcal{L}_{CV}\) da validação \textit{walk-forward}. Assim, a escolha de \(\lambda^{*}\) pode usar várias origens temporais dentro do desenvolvimento, sem consultar uma avaliação final reservada.

No Optuna, cada tentativa, ou \textit{trial}, escolhe um ponto $\lambda$ do espaço de busca, treina o modelo com esse ponto e retorna uma métrica objetivo. O conjunto de tentativas forma um estudo (\textit{study}). A cada tentativa, o amostrador usa o histórico de configurações já avaliadas para propor novos hiperparâmetros. No caso padrão de otimização mono-objetivo, o Optuna usa o \texttt{TPESampler} como amostrador padrão \cite{optunaDocsEfficient2026}. O TPE, ou \textit{Tree-structured Parzen Estimator}, separa os valores de hiperparâmetros associados aos melhores resultados dos demais, ajusta modelos probabilísticos para esses dois grupos e tende a escolher novos valores que maximizem a razão entre regiões promissoras e não promissoras \cite{optunaDocsTPE2026}.

A poda, ou \textit{pruning}, é uma forma de parada antecipada no nível da HPO. Em vez de treinar todas as configurações até o orçamento máximo, o treinamento reporta valores intermediários de validação; o mecanismo de poda decide se uma tentativa deve continuar ou ser interrompida. O \texttt{HyperbandPruner} combina múltiplas instâncias de \textit{Successive Halving}, chamadas de chaves (\textit{brackets}), para distribuir orçamento entre muitas configurações baratas e poucas configurações que recebem mais épocas \cite{li2018hyperband,optunaDocsHyperband2026}. A própria documentação do Optuna apresenta o \texttt{HyperbandPruner} como uma combinação recomendada com \texttt{TPESampler} em resultados de referência para tarefas mono-objetivo \cite{optunaDocsEfficient2026}. Mesmo assim, HPO com orçamento limitado não deve ser interpretada como garantia de ótimo global.

\section{Métricas de Avaliação}
\label{sec:metricas_avaliacao}

A avaliação de modelos de regressão multi-horizonte deve combinar métricas globais e métricas por horizonte. Métricas como erro médio absoluto (\textit{Mean Absolute Error}, MAE), raiz do erro quadrático médio (\textit{Root Mean Squared Error}, RMSE) e erro percentual médio absoluto (\textit{Mean Absolute Percentage Error}, MAPE) são amplamente usadas para avaliar desempenho preditivo \cite{hyndman2021forecasting}. O MAE mede a magnitude média dos erros e é menos sensível a valores extremos que o erro quadrático:

\[
    \text{MAE}(y, \hat{y}) =
    \frac{1}{N}\sum_{i=1}^{N}|y_i - \hat{y}_i|.
\]

A raiz do erro quadrático médio (RMSE) penaliza erros grandes com mais intensidade:

\[
    \text{RMSE}(y, \hat{y}) =
    \sqrt{\frac{1}{N}\sum_{i=1}^{N}(y_i - \hat{y}_i)^2}.
\]

O erro percentual médio absoluto (MAPE) expressa o erro relativo em percentual, com cuidado adicional para valores próximos de zero:

\[
    \text{MAPE}(y, \hat{y}) =
    \frac{100}{N}\sum_{i=1}^{N}\left|\frac{y_i-\hat{y}_i}{|y_i|+\epsilon}\right|.
\]

O coeficiente de determinação ($R^2$) compara o erro do modelo com a variância do alvo:

\[
    R^2(y, \hat{y}) =
    1 - \frac{\sum_{i=1}^{N}(y_i - \hat{y}_i)^2}{\sum_{i=1}^{N}(y_i - \bar{y})^2}.
\]

Quando há valores imputados no alvo, as métricas devem ser calculadas apenas sobre os pontos originalmente observados, usando a mesma lógica da máscara \(m_i\). Para o MAE observado, por exemplo:

\[
    \text{MAE}_{obs}(y, \hat{y}) =
    \frac{\sum_{i=1}^{N}m_i|y_i-\hat{y}_i|}
    {\sum_{i=1}^{N}m_i}.
\]

O mesmo princípio se aplica a RMSE, MAPE e métricas por horizonte. No caso do \(R^2\), a média \(\bar{y}\) e as somas do numerador e do denominador devem ser calculadas no mesmo subconjunto observado. Isso evita tratar valores imputados como verdade de referência e mantém a avaliação vinculada às medições efetivamente observadas \cite{rubin1976inference,cao2018brits,du2023saits}.

Além dessas métricas globais, análises por horizonte, viés médio, percentis altos previstos e erro em eventos de maior concentração podem complementar a avaliação. Essa combinação é necessária porque um modelo pode vencer em erro médio e ainda produzir curvas excessivamente suavizadas, subestimando picos relevantes para a interpretação ambiental.

\clearpage

\section{Resumo do Capítulo}

Este capítulo definiu a base teórica do estudo: PM2,5 como série temporal ambiental, componentes temporais, autocorrelação, previsão multi-horizonte, tratamento de dados faltantes sem vazamento, janelas supervisionadas, variáveis exógenas disponíveis no momento da previsão, redes neurais artificiais, aprendizado profundo, LSTM, Seq2Seq com atenção, \textit{scheduled sampling}, XGBoost, HPO e métricas. A formulação matemática deixa explícita a distinção entre arquiteturas recorrentes, codificador-decodificador e modelos baseados em árvores, preparando a metodologia comparativa dos capítulos seguintes.
